\documentclass[UTF8,fontset=windows,12pt]{ctexart}
%\setcounter{tocdepth}{1}
%\setcounter{secnumdepth}{4}
%\usepackage{ctex} %若要输入中文请取消注释或者把article改成ctexart
% This first part of the file is called the PREAMBLE. It includes
% customizations and command definitions. The preamble is everything
% between \documentclass and \begin{document}.
%\usepackage[OT1]{fontenc}
%\usepackage{fontspec}
%\setmainfont{Times New Roman}
%\setCJKmainfont{Microsoft YaHei}
\usepackage[left=1.91cm,right=1.91cm,top=2.54cm,bottom=2.54cm]{geometry}
\usepackage{amsfonts}
\usepackage{amsmath}
\usepackage{amssymb}
\usepackage[upint,lcgreekalpha]{stix}
\usepackage[type1]{garamondlibre}
%\usepackage{esint}
\usepackage{pgfplots}
\pgfplotsset{compat=1.18}
\usepackage{tikz}
\usetikzlibrary{arrows}
\usepackage{extarrows}
\usepackage[only,llbracket,rrbracket]{stmaryrd}
\usepackage{titlesec}
\usepackage{bm}
\usepackage{graphicx}
\usepackage{appendix}
\usepackage{tikz-cd}
\usepackage{tikz-3dplot}
\usetikzlibrary{arrows.meta, decorations.pathreplacing, calc, patterns}
\usepackage{float}
\usepackage{mathtools}
\usepackage{amsthm}
\usepackage{verbatim}              % better theorem environments
\usepackage{setspace}
\usepackage{enumitem}
\setenumerate[1]{itemsep=0pt,partopsep=0pt,parsep=\parskip,topsep=5pt}
\setitemize[1]{itemsep=0pt,partopsep=0pt,parsep=\parskip,topsep=5pt}
\setdescription{itemsep=0pt,partopsep=0pt,parsep=\parskip,topsep=5pt}
\usepackage{xcolor}
\usepackage[colorlinks,linkcolor=black,anchorcolor=black,citecolor=black]{hyperref}
%\usepackage{apacite}
\usepackage[numbers]{natbib}
%\usepackage{showkeys}
\usepackage{cases}
\usepackage{fancyhdr}
\usepackage[scaled=0.92]{helvet}	% set Helvetica as the sans-serif font
%\usepackage{upgreek}
%\renewcommand{\rmdefault}{ptm}		% set Times as the default text font
%\usepackage{newtxtext}
\bibliographystyle{plain}
% If AMS-LaTeX is used, it can be loaded before or after mtpro2		
% The following loads metro and defines some common MTPro options [2, 4]
%\usepackage[lite,subscriptcorrection,nofontinfo]{mtpro2}
%\pdfmapfile{+mtpro2.map}

%\usepackage{txfonts}

% various theorems, numbered by section

\makeatletter
\renewenvironment{proof}[1][\proofname]{\par
  \pushQED{\qed}%
  \normalfont \topsep6\p@\@plus6\p@\relax
  \trivlist
  \item[\hskip\labelsep
        \bfseries % 关键修改：斜体改为加粗
        #1\@addpunct{.}]\ignorespaces
}{%
  \popQED\endtrivlist\@endpefalse
}
\makeatother
\newtheoremstyle{kaiti-theorem}   % 样式名称
  {3pt}                           % 上方间距
  {3pt}                           % 下方间距
  {\kaishu\upshape}               % 正文字体：中文楷体 + 英文直立
  {}                              % 缩进量
  {\bfseries\upshape}      % 头部字体（定理名）：中文楷体粗体 + 英文直立粗体
  {.}                             % 定理名后标点
  {0.5em}                         % 定理名后间距
  {}                              % 头部说明（留空）

% 应用自定义样式
\theoremstyle{kaiti-theorem}
\newtheorem{thm}{定理}[section]
\newtheorem{nota}[thm]{记号}
\newtheorem{question}{问题}
\newtheorem*{questionn}{思考}
\newtheorem{exe}{习题}[section]
\newtheorem{assump}[thm]{假设}
\newtheorem*{thmm}{定理}
\newtheorem{defn}{定义}[section]
\newtheorem{lem}[thm]{引理}
\newtheorem*{lemm}{引理}
\newtheorem{prop}[thm]{命题}
\newtheorem*{propp}{命题}
\newtheorem*{claim}{断言}
\newtheorem{cor}[thm]{推论}
\newtheorem{conj}[thm]{猜想}
\newtheorem{rmk}{注记}[section]
\newtheorem{ex}{例}[section]
\newtheorem{pf}{证明}
\newtheorem{problem}{作业题}
\numberwithin{equation}{section}





\DeclareMathOperator{\id}{id}
\DeclareMathOperator*{\esssup}{ess\,sup}
\renewcommand{\Re}{\textbf{Re}\,}  
\renewcommand{\Im}{\textbf{Im}\,}  
\newcommand{\R}{\mathbb{R}}  
\newcommand{\FH}{\mathbb{H}}  
\newcommand{\C}{\mathbb{C}}  
\newcommand{\Z}{\mathbb{Z}}
\newcommand{\X}{\mathbb{X}}
\newcommand{\N}{\mathbb{N}}
\newcommand{\Q}{\mathbb{Q}}
\newcommand{\T}{\mathbb{T}}
\newcommand{\TT}{\mathcal{T}}
\newcommand{\TP}{\overline{\partial}{}}
\newcommand{\TJ}{\langle\TP\rangle}
\newcommand{\TL}{\overline{\Delta}{}}
\newcommand{\curl}{\text{curl }}
\newcommand{\dive}{\text{div }}
\newcommand{\bd}[1]{\mathbf{#1}}  % for bolding symbols
\newcommand{\bds}[1]{\boldsymbol{#1}}  % for bolding symbols
\newcommand{\RR}{\mathcal{R}}      % for Real numbers
\newcommand{\sh}{\mathcal{S}}  
\newcommand{\sph}{\mathbb{S}}  
\newcommand{\Om}{\Omega}
\newcommand{\OmT}{{\Omega_T}}
\newcommand{\ut}{{U_T}}
\newcommand{\vp}{\varphi}
\newcommand{\Th}{\Theta}
\newcommand{\Lam}{\Lambda}
\newcommand{\Omb}{{\overline{\Omega}}}
\newcommand{\OmbT}{{\overline{\Omega_T}}}
\newcommand{\ub}{{\overline{U}}}
\newcommand{\ubt}{{\overline{U_T}}}
\newcommand{\q}{\quad}
\newcommand{\p}{\partial}
\newcommand{\pb}{\boldsymbol{\partial}}
\newcommand{\lap}{\Delta}
\newcommand{\dd}{\mathfrak{D}}
\newcommand{\DD}{\mathcal{D}}
\newcommand{\den}{{\bar\rho}_{0}}
\newcommand{\Dt}{{\p}_{t}+v^{k}{\p}_{k}}
\newcommand{\col}[1]{\left[\begin{matrix} #1 \end{matrix} \right]}
\newcommand{\noi}{\noindent}
\newcommand{\nab}{\nabla}
\newcommand{\cnab}{\overline{\nab}}
\newcommand{\cp}{\overline{\partial}{}}
\newcommand{\dx}{\,\mathrm{d}x}
\newcommand{\dxx}{\,\mathrm{d}\bds{x}}
\newcommand{\xxi}{\boldsymbol{\xi}}
\newcommand{\eeta}{\boldsymbol{\eta}}
\newcommand{\dxxi}{\,\mathrm{d}\boldsymbol{\xi}}
\newcommand{\deeta}{\,\mathrm{d}\boldsymbol{\eta}}
\newcommand{\dxi}{\,\mathrm{d}\xi}
\newcommand{\deta}{\,\mathrm{d}\eta}
\newcommand{\dy}{\,\mathrm{d}y}
\newcommand{\dyy}{\,\mathrm{d}\bds{y}}
\newcommand{\dz}{\,\mathrm{d}z}
\newcommand{\dzz}{\,\mathrm{d}\bds{z}}
\newcommand{\dph}{\,\mathrm{d}\varphi}
\newcommand{\dt}{\,\mathrm{d}t}
\newcommand{\dr}{\,\mathrm{d}r}
\newcommand{\dtt}{\,\mathrm{d}\tau}
\newcommand{\dS}{\,\mathrm{d}S}
\newcommand{\dss}{\,\mathrm{d}\sigma}
\newcommand{\dmu}{\,\mathrm{d}\mu}
\newcommand{\dnu}{\,\mathrm{d}\nu}
\newcommand{\dV}{\,\mathrm{d}V}
\newcommand{\ds}{\,\mathrm{d}s}
\newcommand{\drr}{\,\mathrm{d}\rho}
\newcommand{\dist}{\text{dist }}
\newcommand{\vol}{\text{vol}\,}
\newcommand{\volo}{\text{vol}(\Omega)}
\newcommand{\spt}{\text{Spt}\,}
\newcommand{\Tr}{\text{Tr}\,}
\newcommand{\lee}{\langle}
\newcommand{\ree}{\rangle}
\newcommand{\xee}{\langle\xi\rangle}
\newcommand{\xxee}{\langle\boldsymbol{\xi}\rangle}
\newcommand{\xeta}{\langle\eta\rangle}
\newcommand{\xeeta}{\langle\boldsymbol{\eta}\rangle}
\newcommand{\kk}{\kappa}
\newcommand{\lam}{\lambda}
\newcommand{\io}{\int_{\Omega}}
\newcommand{\iu}{\int_{U}}
\newcommand{\is}{\int_{\Sigma}}
\newcommand{\ipo}{\int_{\partial\Omega}}
\newcommand{\ipu}{\int_{\partial U}}
\newcommand{\ig}{\int_{\Gamma}}
\newcommand{\ir}{\int_{\R}}
\newcommand{\ird}{\int_{\R^d}}
\newcommand{\irn}{\int_{\R^n}}
\newcommand{\fpi}{\frac{1}{(2\pi)^{\frac{d}{2}}}}
\newcommand{\ddt}{\frac{\mathrm{d}}{\mathrm{d}t}}
\newcommand{\ddx}{\frac{\mathrm{d}}{\mathrm{d}x}}
\newcommand{\eps}{\varepsilon}
\providecommand{\jump}[1]{\left\llbracket #1 \right\rrbracket }
\providecommand{\len}[1]{\lee #1 \ree }
\providecommand{\ino}[1]{\left\| #1 \right\| }
\providecommand{\bno}[1]{\left| #1 \right| }
\newcommand{\red}{\textcolor{red}}
\newcommand{\green}{\textcolor{green}}
\newcommand{\blue}{\textcolor{blue}}
\newcommand{\nnr}{{[n+1]}}
\newcommand{\nnn}{{[n]}}
\newcommand{\nnl}{{[n-1]}}
\newcommand{\nnll}{{[n-2]}}	
\newcommand{\mmr}{{[m+1]}}
\newcommand{\mmm}{{[m]}}
\newcommand{\mml}{{[m-1]}}
\newcommand{\mmll}{{[m-2]}}

%---------------Special variables--------------
\newcommand{\CC}{\mathfrak{C}}  
\newcommand{\NN}{\mathbf{N}}
\newcommand{\LH}{\mathcal{L}}
\newcommand{\HH}{\mathcal{H}}  
\newcommand{\EE}{\mathcal{E}}  
\newcommand{\FT}{\mathcal{F}}  
\newcommand{\fT}{\mathbf{T}}  
\newcommand{\FA}{\mathbf{A}}  
\newcommand{\MH}{\mathcal{M}}
\newcommand{\NH}{\mathcal{N}}
\newcommand{\PH}{\mathcal{P}}
\newcommand{\VH}{\mathcal{V}}
\newcommand{\XX}{\mathfrak{X}}
\newcommand{\xx}{{\bds{x}}}
\newcommand{\bx}{\mathbf{x}}
\newcommand{\by}{\mathbf{y}}
\newcommand{\bp}{\mathbf{p}}
\newcommand{\bq}{\mathbf{q}}
\newcommand{\pp}{{\bds{p}}}
\newcommand{\qq}{{\bds{q}}}
\newcommand{\uu}{\mathbf{u}}
\newcommand{\ww}{\mathbf{w}}
\newcommand{\vv}{\mathbf{v}}
\newcommand{\yy}{{\bds{y}}}
\newcommand{\zz}{{\bds{z}}}
\newcommand{\zo}{\mathbf{0}}
\newcommand{\ee}{\mathbf{e}}
\newcommand{\fa}{\mathbf{a}}
\newcommand{\fb}{\mathbf{b}}
\newcommand{\fc}{\mathbf{c}}
\newcommand{\ff}{\bds{f}}
\newcommand{\fr}{\mathbf{r}}
\newcommand{\fh}{\mathbf{h}}
\newcommand{\fm}{\mathbf{m}}
\newcommand{\fg}{\mathbf{g}}
\newcommand{\FF}{\mathbf{F}}
\newcommand{\GG}{\mathbf{G}}
\newcommand{\BB}{\mathbf{B}}



\begin{document}
\title{\LARGE{\bf  2026年春季学期现代偏微分方程作业一}}
\author{Sobolev空间}
\date{截止时间：2026.3.24下课前（电子版截至当天23:59:59）}

\maketitle

第1-5题为Sobolev函数光滑逼近相关，第6-11题为Sobolev嵌入定理相关。其中第5、8题选做。

\begin{problem}
设$f\in H^1(\R^d)$. 定义光滑逼近 $f_\eps:=\eta_\eps*f$，其中 $\eta_\eps$ 是附录B.2定义的磨光核（亦见 Evans PDE 附录 C.5 节）。证明如下结论：
\begin{itemize}
\item [(1)] $\|f_\eps\|_{H^1}\leqslant \|f\|_{H^1}$, 存在不依赖$\eps$的常数 $C>0$ 使得 $\|f_\eps\|_{H^{2}}\leqslant C\eps^{-1}\|f\|_{H^1}$.
\item [(2)] 存在不依赖$\eps$的常数 $C>0$ 使得 $\|f_\eps - f\|_{L^2}\leqslant C\eps\|f\|_{H^1}$.
\item [(3)] 当 $\eps\to 0_+$ 时，是否一定成立 $\eps^{-1}\|f_\eps - f\|_{L^2}\to 0$？证明你的结论。
\item [(4)] $C_c^\infty(\R^d)$ 在 $H^1(\R^d)$ 中是稠密的。
\end{itemize}
提示：(3) 注意磨光核 $\eta$ 本身是径向对称函数，(4) 先作截断再磨光。
\end{problem}


\begin{problem}[习题1.2.3] 证明讲义上的命题1.2.4(2). 即设$1\leqslant p<\infty$, $f\in W^{1,p}(U)$, $F\in C^1(\R)$, $F'\in L^{\infty}(\R)$, $F(0)=0$, 则 $F(f)\in W^{1,p}(U)$, 且 $\p_i(F(f))=F'(f)\p_i f$对$i=1,\cdots,d$在 $U$ 中几乎处处成立。若$U$在$\R^d$中的Lebesgue测度有限，则$F(0)=0$不是必要的。
\end{problem}

\begin{problem}[习题1.2.4]
证明$\|\nab u\|_{L^2(U)}^2\leqslant C\|u\|_{L^2(U)}\|\nab^2 u\|_{L^2(U)}$  对任何 $u\in C_c^{\infty}(U)$ 成立，其中 $\|\nab^2 u\|_{L^2(U)}^2=\iu\sum\limits_{i,j=1}^d(\p_i\p_j u)^2\dxx$. 然后推广到 $u\in H_0^1(U)\cap H^2(U)$，这里需假设 $U$ 有界且 $\p U$ 光滑。
\end{problem}

\begin{problem}[习题1.2.6]
设 $U\subset\R^d$ 是\textbf{连通}开集且 $f\in W^{1,p}(U)$ 满足 $\nab f=\bd{0}$ 在 $U$ 中几乎处处成立。证明：$f$ 几乎处处等于一个常值函数。

提示：本题\textbf{不能使用Poincar\'e不等式}来证明，因为该题结论在Poincar\'e不等式的证明中要用到。 在 $V\Subset U$ （比如说可以取成一个球） 中考虑 $f_\eps:=\eta_\eps*f$，并证明对充分小的 $\eps>0$，$f_\eps=C_\eps$ 在 $V$ 中几乎处处成立。然后使用定理 C.2.1 来证明 $\|f_\eps\|_{L^p}$ 关于 $\eps$ 一致有界，因此 $\{C_\eps\}$ 也一致有界，进而$\{f_\eps\}$ 的一个子列收敛到某个常数$C$。最后证明 $\{ \xx \in U : \exists r > 0 \text{ 使得 } B(\xx, r) \subset U \text{ 且 } f = C\text{ 几乎处处在 } B(\xx, r) \text{ 中成立} \}$ 在 $U$ 中既开又闭。
\end{problem}

\begin{problem}[习题1.3.3，选做]
设$f\in H_0^1(\R_+^d)\cap H^2(\R_+^d)$. 证明：对任意$1\leqslant i\leqslant d-1$, 偏导数$\p_{x_i} f$也属于$H_0^1(\R_+^d)$.
\end{problem}


\begin{problem}[习题1.4.3]
设 $U\subset\R^d$ 是有界\underline{区域}，函数 $u\in H^1(U)$ 满足：存在$\alpha\in(0,1)$使得集合 $Z:=\{\xx\in U|u(\xx)=0\}$ 满足 $\LH^d(Z)\geqslant \alpha\LH^d(U)$. 证明：存在仅依赖于 $d,\alpha$ 的常数 $C>0$，使得\[
\iu u^2\dxx\leqslant C\iu |\nab u|^2\dxx.
\]

提示：在 $U\backslash Z$ 中，将 $u^2$ 写为 $(u-(u)_U+(u)_U)^2$ 并对 $(u-(u)_U)^2$ 应用 Poincaré 不等式。$(u)_U^2$ 的贡献将被所需不等式的左边吸收，这是因为 $U\backslash Z$ 的测度严格小于 $U$ 的测度。\textbf{我在知乎和个人主页上写的答案里这道题的证明有点问题，请不要去抄。}
\end{problem}


\begin{problem}[习题1.4.7]
设 $U\subset\R^d$ 是边界光滑的有界区域。证明：$H^2(U)$ 紧嵌入到 $H^1(U)$，且对任意 $\eps>0$存在常数 $C_\eps >0$ 使得
 $$\|\nab u\|_{L^2(U)}\leqslant \eps\|u\|_{H^2(U)}+C_\eps \|u\|_{L^2(U)},\quad \forall u\in H^2(U).$$
\end{problem}

\begin{problem}[问题1.4.1，选做]
设$U\subset\R^d$ 是边界光滑的有界区域，开球$B\Subset U$. 对$\eps\in(0,1)$设$u_\eps$是如下方程的光滑解
\begin{equation}
\begin{cases}
-\lap u_\eps + \eps^{-1}(u_\eps -f)\mathbf{1}_{B}=0&\text{ in }U,\\
u_\eps=0&\text{ on }\p U,
\end{cases}
\end{equation}其中 $f\in H_0^1(U)$ 是给定的函数。证明：$\|\nab u_\eps\|_{L^2(U)}$ 关于 $\eps$ 一致有界，且 $u_\eps \xrightarrow{L^2(B)} f$.
\end{problem}

\begin{problem}[问题1.4.4，Strauss径向引理]
设$d\geqslant 2$, $u\in H_{\text{rad}}^1(\R^d)$ (即$u\in H^1(\R^d)$是径向函数，$u(\xx)$取值只依赖$r:=|\xx|$). 本题可默认$C_c^\infty(\R^d)$ 在 $H_{\text{rad}}^1(\R^d)$ 中是稠密的，进而你可以直接对 $C_c^\infty$ 径向函数证明。
\begin{itemize}
	\item [(1)] 利用微积分基本定理和$u(\infty)=0$证明：$|u(r)|^2\leqslant 2r^{1-d}\int_r^{\infty} |u(s)||u'(s)|s^{d-1}\ds$.
	\item [(2)] 证明：$|u(r)|\leqslant Cr^{-(d-1)/2}\|u\|_{L^2(\R^d)}^{1/2}\|\nab u\|_{L^2(\R^d)}^{1/2}$.
\end{itemize}

提示：对(1)右边的积分用Cauchy-Schwarz不等式，再用积分的极坐标表示凑出 $u$ 和 $\nab u$ 的 $L^2$ 范数。
\end{problem}

\begin{problem}[问题1.4.5]
设$d\geqslant 2$, $2<q<2^*:=\frac{2d}{d-2}$, 证明：$H_{\text{rad}}^1(\R^d)\hookrightarrow L^q(\R^d)$ 是紧嵌入。

提示：在$|\xx|=R$处作截断，当$|\xx|>R$时用Strauss径向引理，当$|\xx|\leqslant R$时已经有紧性。
\end{problem}
\begin{problem}[习题1.5.1]
设$U\subset\R^d$ 是边界Lipschitz 的有界开集，正整数$k>\frac{d}{2}$.  证明：存在依赖于 $k,d,U$ 的常数 $C>0$，使得对任意$f,g\in H^k(U)$成立如下不等式
\begin{align*}
\|fg\|_{H^k(U)}\leqslant C\|f\|_{H^k(U)}\|g\|_{H^k(U)}.
\end{align*}
\end{problem}

\end{document}
